\documentclass[tikz,border=10pt]{standalone}
\usepackage{ctex}
\usepackage{tikz}

% ===== 自定义金色 =====
\definecolor{gold}{RGB}{255, 215, 0}
\definecolor{darkgold}{RGB}{184, 134, 11}

% ===== 手绘五角星宏（修复坐标解析问题）=====
% #1=x, #2=y, #3=半径, #4=填充色, #5=描边色
\newcommand{\drawstar}[5]{%
	\begin{scope}[shift={({#1},{#2})}]
		\fill[#4, draw=#5, line width=0.5pt]
		(90:#3)
		\foreach \a in {162,234,306,18} {
			-- (\a:#3)
			-- (\a+36:0.382*#3)
		} -- cycle;
	\end{scope}
}

\begin{document}

	\begin{tikzpicture}[scale=1]

		% ============================================================
		% ===== 背景层 =====
		% ============================================================
		\fill[red!90!black] (0,0) rectangle (16,24);

		\fill[red!50!black, opacity=0.55] (8,13) circle (13);
		\fill[red!60!black, opacity=0.45] (8,13) circle (10);
		\fill[red!70!black, opacity=0.35] (8,13) circle (7);

		\foreach \i in {0,0.5,...,6} {
			\fill[red!75!black, opacity={0.05+0.06*\i}]
			(0,\i) rectangle (16,\i+0.5);
		}

		% ============================================================
		% ===== 放射光芒 =====
		% ============================================================
		\foreach \angle in {0,15,...,345} {
			\draw[gold!25!red, line width=0.6pt, opacity=0.35]
			(8,13) -- ++(\angle:18);
		}
		\foreach \angle in {7.5,22.5,...,352.5} {
			\draw[gold!40!red, line width=1.2pt, opacity=0.18]
			(8,13) -- ++(\angle:16);
		}

		% ============================================================
		% ===== 多层华丽边框 =====
		% ============================================================
		\draw[gold, line width=4pt] (0.45,0.45) rectangle (15.55,23.55);
		\draw[gold!85!white, line width=1pt] (0.75,0.75) rectangle (15.25,23.25);
		\draw[gold!70!white, line width=0.6pt, dashed, dash pattern=on 4pt off 3pt]
		(1.05,1.05) rectangle (14.95,22.95);

		\foreach \x/\y in {0.75/0.75, 15.25/0.75, 0.75/23.25, 15.25/23.25} {
			\fill[gold] (\x,\y) circle (0.18);
			\draw[gold!80!white, line width=0.6pt] (\x,\y) circle (0.35);
			\draw[gold!60!white, line width=0.4pt] (\x,\y) circle (0.55);
		}

		\foreach \y in {0.75, 23.25} {
			\fill[gold] (8,\y) circle (0.18);
			\draw[gold!80!white, line width=0.6pt] (8,\y) circle (0.35);
			\drawstar{8}{\y}{0.3}{gold!90!white}{darkgold!80!black}
		}

		% ============================================================
		% ===== 顶部卷轴飘带 =====
		% ============================================================
		\fill[black, opacity=0.35]
		(3.3,21.2) -- (12.9,21.2) -- (12.5,20.4) -- (3.7,20.4) -- cycle;
		\fill[gold!90!white]
		(3.2,21.3) -- (12.8,21.3) -- (12.4,20.5) -- (3.6,20.5) -- cycle;
		\fill[gold!70!white]
		(3.2,21.3) -- (2.4,20.9) -- (3.6,20.5) -- cycle;
		\fill[gold!70!white]
		(12.8,21.3) -- (13.6,20.9) -- (12.4,20.5) -- cycle;
		\node[font=\bfseries\zihao{3}, text=red!80!black] at (8,20.9)
		{中国人民解放军建军};

		% ============================================================
		% ===== 中心军徽 =====
		% ============================================================
		\draw[gold, line width=2.5pt] (8,15.5) circle (3.6);
		\draw[gold!80!white, line width=0.8pt] (8,15.5) circle (3.95);
		\draw[gold!60!white, line width=0.5pt, dashed, dash pattern=on 3pt off 2pt]
		(8,15.5) circle (4.2);

		\drawstar{8.06}{15.44}{2.1}{black!50}{black!50}
		\drawstar{8}{15.5}{2.1}{gold}{darkgold}

		\fill[red!85!black] (8,15.5) circle (1.05);

		\node[font=\bfseries\zihao{2}, text=gold] at (8,15.5) {八一};

		% 环绕五颗小星（用 \pgfmathsetmacro 先算好坐标）
		\foreach \angle in {90,162,234,306,18} {
			\pgfmathsetmacro{\sx}{8+3.0*cos(\angle)}
			\pgfmathsetmacro{\sy}{15.5+3.0*sin(\angle)}
			\drawstar{\sx}{\sy}{0.5}{gold!90!white}{darkgold!80!black}
		}

		% ============================================================
		% ===== 主标题区 =====
		% ============================================================
		\node[font=\bfseries\zihao{0}, text=black!55, opacity=0.6,
		xshift=2.5pt, yshift=-2.5pt] at (8,11.2) {建军百年};
		\node[font=\bfseries\zihao{0}, text=gold] at (8,11.2) {建军百年};

		\node[font=\bfseries\zihao{2}, text=gold!95!white] at (8,9.6)
		{1927 \; — \; 2027};

		% ============================================================
		% ===== 分隔装饰线 =====
		% ============================================================
		\draw[gold, line width=2pt] (3.5,8.8) -- (12.5,8.8);
		\draw[gold!80!white, line width=0.8pt] (4.5,8.55) -- (11.5,8.55);
		\drawstar{8}{8.8}{0.35}{gold!90!white}{darkgold!80!black}

		% ============================================================
		% ===== 副标题 =====
		% ============================================================
		\node[font=\bfseries\zihao{3}, text=gold!95!white] at (8,7.6)
		{峥嵘岁月 \; · \; 铁血军魂};
		\node[font=\zihao{4}, text=gold!85!white] at (8,6.8)
		{传承红色基因 \; · \; 担当强军重任};

		% ============================================================
		% ===== 底部三大标语 =====
		% ============================================================
		\foreach \x/\txt in {3.5/听党指挥, 8/能打胜仗, 12.5/作风优良} {
			\draw[gold!80!white, line width=1pt, rounded corners=3pt]
			(\x-1.6,4.3) rectangle (\x+1.6,5.5);
			\node[font=\bfseries\zihao{4}, text=gold] at (\x,4.9) {\txt};
		}

		% ============================================================
		% ===== 底部年份与落款 =====
		% ============================================================
		\node[font=\bfseries\zihao{3}, text=gold!90!white] at (8,3.3)
		{1927 \; — \; 2027};

		\node[font=\zihao{5}, text=gold!60!white] at (8,1.8)
		{纪念中国人民解放军建军一百周年};

		% ============================================================
		% ===== 底部小星装饰 =====
		% ============================================================
		\foreach \x in {4.2,5.4,6.6,9.4,10.6,11.8} {
			\drawstar{\x}{1.1}{0.2}{gold!90!white}{darkgold!80!black}
		}

		% ============================================================
		% ===== 两侧麦穗勋章装饰 =====
		% ============================================================
		\begin{scope}[shift={(1.6,10)}, rotate=0]
			\draw[gold!85!white, line width=1.8pt]
			(0,5) .. controls (-0.6,3) and (-0.4,1) .. (0,-1);
			\foreach \y in {4.5, 3.9, 3.3, 2.7, 2.1, 1.5, 0.9, 0.3} {
				\draw[gold!80!white, line width=0.7pt]
				(0,\y) .. controls (-0.5,\y+0.35) and (-0.9,\y+0.25) .. (-1.2,\y+0.6);
				\fill[gold!80!white] (-1.15,\y+0.55) circle (0.07);
			}
		\end{scope}

		\begin{scope}[shift={(14.4,10)}, rotate=180]
			\draw[gold!85!white, line width=1.8pt]
			(0,5) .. controls (-0.6,3) and (-0.4,1) .. (0,-1);
			\foreach \y in {4.5, 3.9, 3.3, 2.7, 2.1, 1.5, 0.9, 0.3} {
				\draw[gold!80!white, line width=0.7pt]
				(0,\y) .. controls (-0.5,\y+0.35) and (-0.9,\y+0.25) .. (-1.2,\y+0.6);
				\fill[gold!80!white] (-1.15,\y+0.55) circle (0.07);
			}
		\end{scope}

		% ============================================================
		% ===== 飘带装饰 =====
		% ============================================================
		\draw[gold!70!white, line width=0.8pt, opacity=0.8]
		(2.5,19.5) .. controls (5,19.1) and (11,19.1) .. (13.5,19.5);
		\draw[gold!60!white, line width=0.5pt, opacity=0.7]
		(2.8,19.1) .. controls (5,18.8) and (11,18.8) .. (13.2,19.1);

	\end{tikzpicture}

\end{document}